\title{QLoRA: Efficient Finetuning of Quantized LLMs}

\begin{abstract}
We introduce \textsc{QLoRA}, a finetuning method designed for optimal memory efficiency, allowing a 65B parameter model to be finetuned on a single 48GB GPU without compromising the performance traditionally associated with 16-bit finetuning. The core of \textsc{QLoRA} lies in backpropagating gradients through a static, 4-bit quantized language model into Low Rank Adapters~(LoRA). Our leading model series, termed \textbf{Guanaco}, surpasses all previously released open models on the Vicuna benchmark, achieving 99.3\% of ChatGPT's effectiveness, yet requiring just 24 hours of training on one GPU. To accomplish this, \textsc{QLoRA} leverages several memory-saving innovations that do not detract from performance: (a) the introduction of 4-bit NormalFloat (NF4), a data type theoretically optimal for weights following a normal distribution; (b) Double Quantization, which minimizes memory needs further by quantizing even the quantization constants; and (c) Paged Optimizers, which help control sudden spikes in memory usage. Leveraging \textsc{QLoRA}, we have finetuned more than 1,000 models, conducting an extensive assessment of instruction-following capabilities and chatbot effectiveness over 8 different instruction datasets, multiple architectures (LLaMA, T5), and larger model sizes that are otherwise impractical with standard finetuning techniques (such as models with 33B and 65B parameters). Our findings demonstrate that using QLoRA finetuning on a compact, high-quality training set yields state-of-the-art results, often with smaller models than prior top-performing systems. We provide a thorough evaluation of chatbot performance using both human and GPT-4 assessments, establishing that GPT-4-based evaluations offer a cost-effective and sound alternative to traditional human judgment. Additionally, our work highlights the unreliability of current chatbot benchmarks for accurately assessing chatbot proficiency. Through a targeted failure analysis (“lemon-picked”), we highlight the specific areas in which \textbf{Guanaco} falls short relative to ChatGPT. All models and code, including CUDA kernels for 4-bit training, are made publicly available.\footnote{\url{https://github.com/artidoro/qlora} and \url{https://github.com/TimDettmers/bitsandbytes}}
\end{abstract}

\section{Introduction}

Adapting large language models (LLMs) through finetuning has proven to be an effective strategy for enhancing their capabilities \citep{min2021metaicl, wei2021finetuned, ouyang2022training, wang2022super, wang2022self, liu2022few}, as well as for promoting favorable behaviors or mitigating unfavorable ones \citep{ouyang2022training,askell2021general,bai2022training}. Nevertheless, the computational costs of finetuning models at the largest scales are often prohibitively high; for instance, standard 16-bit finetuning of the LLaMA 65B model~\citep{touvron2023llama} demands upwards of 780 GB of GPU memory. Although quantization techniques have recently been employed to decrease memory consumption for LLMs \citep{dettmers2022llm, dettmers2022case, frantar2022gptq, xiao2022smoothquant}, these advances are confined to inference phases and are not applicable during training \citep{wortsman2023stable}.

In this work, we present the first evidence that a quantized 4-bit model can be finetuned without any sacrifice in performance. Our approach, termed \textsc{QLoRA}, introduces a novel, high-precision quantization scheme that compresses a pretrained model to 4 bits, and augments it with a compact set of trainable Low-rank Adapter weights \citep{hu2021lora} 

that are updated through backpropagation, with gradients flowing through the quantized weights.

By adopting \textsc{QLoRA}, the mean memory overhead for finetuning a 65B parameter model drops from over 780GB to under 48GB of GPU memory, all while maintaining parity in computational speed and predictive accuracy with a baseline fully finetuned at 16-bit precision. This advance dramatically broadens the feasibility of LLM finetuning, making it practical to customize the largest open-source models on a single GPU. Leveraging \textsc{QLoRA}, we develop the \textbf{Guanaco} suite of models: notably, our second-best model attains 97.8\% of ChatGPT’s performance on the Vicuna~\citep{vicuna2023} benchmark and trains in under 12 hours on a standard consumer GPU. With a single professional GPU, we further achieve 99.3\% of ChatGPT’s performance using our largest model after just 24 hours, effectively matching ChatGPT’s score on this evaluation. For deployment, our smallest \textbf{Guanaco} model (7B parameters) operates with only 5 GB of memory and surpasses the 26 GB Alpaca model by more than 20 percentage points in Vicuna benchmark performance (see Table~\ref{tbl:vicuna_eval}).

\textsc{QLoRA} incorporates several novel techniques to minimize memory consumption while maintaining model effectiveness: (1) \textbf{4-bit NormalFloat}, a quantization format that is theoretically optimal for normally distributed variables and empirically surpasses both 4-bit Integer and 4-bit Float types; (2) \textbf{Double Quantization}, a scheme in which quantization constants themselves are quantized, leading to an average memory saving of approximately 0.37 bits per parameter (translating to roughly 3 GB for a 65B parameter model); (3) \textbf{Paged Optimizers}, which leverage NVIDIA’s unified memory to eliminate memory spikes from gradient checkpointing encountered during training on lengthy sequences in mini-batches. Together, these advances yield a refined LoRA methodology that employs adapters at each layer of the model, thereby largely avoiding the compromises in accuracy observed in previous studies.

Owing to \textsc{QLoRA}’s improved memory efficiency, we are able to systematically investigate instruction tuning and chatbot model performance at scales infeasible for traditional finetuning approaches, given their high memory requirements. This work includes the training of over 1,000 distinct models spanning multiple instruction tuning datasets, model families, and parameter sizes from 80M up to 65B. We demonstrate not only that \textsc{QLoRA} attains parity with standard 16-bit finetuning in terms of performance (\S\ref{sec:comp_to_std}) and enables the development of the advanced chatbot \textbf{Guanaco} (\S\ref{sec:pushing}), but also provide detailed analysis of the resulting models. Our findings reveal, firstly, that dataset quality far outweighs dataset scale in determining outcomes—for instance, a curated dataset containing 9k samples (OASST1) surpassed the performance of a substantially larger, 450k-sample set (FLAN v2, subsampled) on chatbot-related tasks, even when both datasets are designed to facilitate instruction-following generalization. Secondly, our results indicate that achieving strong scores on the Massive Multitask Language Understanding (MMLU) benchmark does not necessarily translate to superior performance on the Vicuna chatbot benchmark and vice versa, emphasizing that the relevance of the training data to the specific task is of greater significance than its overall size.

Additionally, we conduct a thorough assessment of chatbot capabilities employing both human judges and GPT-4 for evaluation purposes. Our approach utilizes a tournament-style benchmarking framework, wherein different models are pitted against each other in head-to-head matches to generate the optimal response to a specific prompt. The victor of each matchup is determined either by GPT-4 or by human annotators. We then aggregate the outcomes of these tournaments into Elo scores~\citep{elo1967proposed, elo1978rating}, which are used to establish a performance hierarchy among the chatbots. Our findings indicate substantial agreement between GPT-4 and human evaluators with respect to overall model rankings, though notable discrepancies do occasionally occur. Consequently, we emphasize that although model-based evaluations offer a cost-effective alternative to human assessment, they are accompanied by inherent ambiguities.

To complement our benchmarking results, we perform a qualitative examination of the \textbf{Guanaco} models, discussing examples of successful and unsuccessful outputs that elude purely quantitative metrics.

We make available all model outputs, accompanied by human and GPT-4 annotation data, to support future research. Our full codebase and CUDA kernels are open-sourced, and our methods are incorporated within the Hugging Face transformers framework \citep{wolf2019huggingface} to ensure broad accessibility. Furthermore, we release a suite of adapters for 7/13/33/65B parameter models, each trained on eight distinct instruction-following datasets, culminating in 32 unique open-source, fine-tuned models.

\section{Background}
\paragraph{Block-wise k-bit Quantization} Quantization involves transforming an input from a high-precision representation to a lower-precision one. This typically entails converting data types of greater bit-width to those with fewer bits; for instance, reducing 32-bit floating-point values to 8-bit integers. To effectively map the input data onto the full representational range of the target, low-bit data type, it is common practice to normalize by the absolute maximum value of the input, usually formatted as a tensor. Consider quantizing an FP32 tensor into an Int8 tensor constrained to $[-127, 127]$:

\begin{equation}
    \mathbf{X}^{ \text{Int8}} = \text{round}\left( \frac{127}{{\text{absmax}}(\mathbf{X}^{{\text{FP32}}})}\mathbf{X}^{{\text{FP32}}} \right) = \text{round}(c^{\text{FP32}}\cdot \mathbf{X}^{{\text{FP32}}}),
\end{equation}

where $c$ denotes the {\em quantization constant} or {\em quantization scale}. The process can be reversed via dequantization as follows:

\begin{equation}
    \text{dequant}(c^{\text{FP32}}, \mathbf{X}^{\text{Int8}}) = \frac{\mathbf{X}^{\text{Int8}}}{c^{\text{FP32}}} = \mathbf{X}^{\text{FP32}}
\end{equation}

A limitation of this method arises when the input tensor contains values with very large magnitude (outliers); in such scenarios, the quantization bins—or specific bit patterns—are inefficiently assigned, with some bins corresponding to little or no quantized values. To address the problem of outliers, a standard solution is to partition the input tensor into several blocks, each of which is quantized separately using its distinct quantization constant $c$. More precisely, the input tensor $\mathbf{X}\in \mathbb{R}^{b\times h}$ is divided into $n$ contiguous blocks of equal size $B$. This is accomplished by flattening the tensor and slicing the sequence into ${n = ({b\times h} )/{B} }$ segments. Each block is then quantized independently according to Equation~1, yielding a quantized tensor together with $n$ unique quantization constants $c_i$.

\paragraph{Low-rank Adapters} The LoRA (Low-rank Adapter) finetuning approach \citep{hu2021lora} reduces memory usage by introducing a set of lightweight, trainable adapter parameters, while keeping the main model weights unchanged throughout training. During stochastic gradient descent, gradients propagate through the frozen pretrained weights to update only the adapter, thereby driving loss minimization. LoRA modifies the standard linear projection by incorporating an additional low-rank, factorized projection. For a given projection ${\mathbf{X} \mathbf{W} = \mathbf{Y}}$, where $\mathbf{X}\in  \mathbb{R}^{b\times h}$ and $\mathbf{W}\in  \mathbb{R}^{h\times o}$, LoRA applies:

\begin{equation}
    \mathbf{Y} = \mathbf{X}\mathbf{W} + s\mathbf{X}\mathbf{L}_1\mathbf{L}_2,
\end{equation}

in which $\mathbf{L}_1\in  \mathbb{R}^{h\times r}$, $\mathbf{L}_2\in  \mathbb{R}^{r\times o}$, and $s$ denotes a scaling factor.

\paragraph{Memory Requirement of Parameter-Efficient Finetuning} The consideration of memory usage in LoRA training, specifically concerning both the number and capacity of adapters, is a critical aspect. Due to the extremely low memory overhead introduced by LoRA, incorporating additional adapters to enhance model performance is feasible without causing substantial growth in overall memory consumption. Although LoRA is a Parameter Efficient Finetuning (PEFT) approach, the dominant contributor to memory consumption during large language model finetuning is actually the storage of activation gradients rather than the LoRA parameters themselves. For instance, when finetuning a 7B LLaMA model on FLAN v2 with a batch size of 1 and utilizing LoRA weights amounting to the common 0.2\% of the base model weights~\citep{hu2021lora,liu2022few}, the memory required for LoRA input gradients is approximately 567 MB, whereas the memory occupied by the LoRA parameters themselves is just 26 MB. Implementing gradient checkpointing~\citep{chen2016training} further lowers the average memory demand for input gradients to 18 MB per sequence, which still surpasses the total LoRA parameter memory. In stark contrast, the main 4-bit model consumes 5,048 MB. These observations underscore the importance of gradient checkpointing and demonstrate that further minimizing the size of LoRA parameters only achieves slight improvements in memory usage. Therefore, it is practical to employ a larger number of adapters without meaningfully elevating the memory required for training (refer to Appendix~\ref{app:memory} for comprehensive details). As addressed in subsequent sections, this flexibility is vital for regaining the performance of models trained in full 16-bit precision.

\section{\textsc{QLoRA} Finetuning}

\textsc{QLoRA} enables effective finetuning in 4-bit precision by leveraging two novel methods: 4-bit NormalFloat (NF4) quantization and Double Quantization. Furthermore, we present Paged Optimizers to mitigate abrupt memory usage increases during gradient checkpointing—an issue that has frequently led to out-of-memory failures when finetuning large models on a single machine.

In \textsc{QLoRA}, model parameters are stored using a single low-precision format—commonly 4-bit—and computations are performed in a separate data type, typically BFloat16. Operationally, whenever a weight tensor is accessed in \textsc{QLoRA}, it is first dequantized to BFloat16, followed by matrix multiplication at 16-bit precision.

We proceed by examining each element constituting \textsc{QLoRA}, and then provide a precise definition of \textsc{QLoRA}.

\paragraph{4-bit NormalFloat Quantization} The NormalFloat (NF) format extends the principles of Quantile Quantization \citep{dettmers2022optimizers}, which is considered theoretically optimal by information theory: it assigns an equal portion of tensor entries to each quantization level. Quantile quantization achieves this by employing the empirical cumulative distribution function to estimate input quantiles.

A primary drawback of quantile quantization lies in the computational cost associated with quantile estimation. As a remedy, efficient quantile approximation algorithms, such as SRAM quantiles~\citep{dettmers2022optimizers}, are employed. However, due to the inherent inaccuracy of these approximations, significant quantization errors may be introduced, especially for outliers, which can be particularly consequential.

Both the high cost of precise quantile computation and the risk of approximation errors can be avoided when the input tensors originate from a distribution that remains consistent except for a quantization scaling factor. Under these circumstances, input tensors share their quantile structure, allowing for exact and tractable quantile computation.

Pretrained neural network weights are typically distributed according to a normal distribution with mean zero and standard deviation $\sigma$ (refer to Appendix~\ref{app:normality}). To unify all weights under a single reference distribution, we can adjust $\sigma$ so that the resulting distribution precisely aligns with the interval permitted by our chosen data type. For our purposes, we define this interval as $[-1, 1]$. Consequently, both the quantile locations corresponding to the data type and the weights themselves require normalization within these bounds.

The data type that is optimal from an information theory standpoint for normal distributions centered at zero with variable standard deviation $\sigma$ (confined to $[-1, 1]$) is determined through the following procedure: (1) identify the $2^k+1$ quantiles of a standard $N(0,1)$ distribution, thereby producing a $k$-bit quantile-quantized data type suitable for Gaussian distributions, (2) map these quantile values to the interval $[-1, 1]$, and (3) apply absolute maximum rescaling to an input weight tensor so its values are also mapped onto $[-1, 1]$ before quantization.

With the value ranges of the weights and data type aligned, conventional quantization methods can be used. The third step essentially rescales the standard deviation of the weights to correspond to the standard deviation inherent to the $k$-bit data type. More rigorously, we define the $2^k$ representative quantile values $q_i$ for the data type by

\begin{equation}
q_i  = \frac{1}{2}\left( Q_X\left(\frac{i}{2^k+1}\right) + Q_X\left(\frac{i+1}{2^k+1}\right)\right),
\end{equation}

where $Q_X(\cdot)$ denotes the quantile function associated with the standard normal distribution $N(0,1)$. One limitation of symmetric k-bit quantization is the lack of an explicit zero representation within the quantized set, which is problematic when precisely encoding zeros for padding or similar elements without incurring quantization error. To guarantee the existence of a discrete zero point at $0$ and simultaneously utilize the entire $2^k$ set of codes for a k-bit representation, we formulate an asymmetric scheme. This involves calculating quantiles $q_i$ over two distinct regions: $2^{k-1}$ quantiles for the negative range, and $2^{k-1}+1$ quantiles for the positive range, merging the resulting sets, and discarding one of the duplicate zeros arising at the boundary. We refer to this scheme—characterized by having uniform expected counts in each bin for the underlying distribution—as the \emph{k-bit NormalFloat} (NFk) data type. NFk achieves information-theoretic optimality for quantizing zero-mean, normally distributed variables. Specific numerical values for this data type can be found in Appendix~\ref{app:nf4}.

\paragraph{Double Quantization{}} 
We propose the technique of \emph{Double Quantization} (DQ), wherein quantization constants themselves are subject to a second quantization stage to further reduce memory consumption. Although precise 4-bit quantization requires a relatively small blocksize~\citep{dettmers2022case}, the associated quantization constants can introduce non-trivial memory overhead. For instance, quantization with 32-bit constants and a blocksize of 64 for $\mathbf{W}$ adds, on average, $32/64 = 0.5$ bits per parameter. Double Quantization is introduced to mitigate the storage requirements associated with these constants.

In detail, Double Quantization entails taking the primary quantization constants $c_2^{\text{FP32}}$ and applying a secondary quantization. This process yields secondary quantized constants $c_2^{\text{FP8}}$ and a new set of quantization constants $c_1^{\text{FP32}}$ for this second stage. For the second quantization step, we employ 8-bit floating-point types and a larger blocksize of 256, given that empirical findings (e.g.,~\citep{dettmers2022case}) indicate no accuracy loss at this resolution. Since $c_2^{\text{FP32}}$ are strictly positive values, we normalize by subtracting the mean, recentering the data and enabling the use of symmetric quantization. Calculating the resultant memory savings: for a blocksize of 64, this strategy decreases per-parameter storage for constants from $0.5$ bits ($32/64$) down to $0.127$ bits ($8/64 + 32/(64 \cdot 256)$), thus saving $0.373$ bits per parameter.

\paragraph{Paged Optimizers} leverage the NVIDIA unified memory~\footnote{\tiny \url{https://docs.nvidia.com/cuda/cuda-c-programming-guide}} capability, which seamlessly manages automatic transfers of memory pages between the CPU and GPU, thus preventing errors in GPU operations even when GPU memory is temporarily exhausted. This mechanism operates in a manner similar to traditional paging systems between CPU RAM and storage drives. In our approach, we allocate the optimizer state memory as paged memory so that when GPU memory limits are reached, these optimizer states are automatically swapped to CPU RAM. Conversely, they are transferred back to the GPU as soon as they are required during optimizer update computations.

\paragraph{\textsc{QLoRA}.} Building upon the elements introduced earlier, we formalize \textsc{QLoRA} for a quantized base model’s single linear layer, enhanced with a single LoRA adapter, as described by:

\begin{equation}
    \mathbf{Y}^{\text{BF16}} =  \mathbf{X}^{\text{BF16}} \text{doubleDequant}(c_1^{\text{FP32}}, c_2^{\text{k-bit}}, \mathbf{W}^{\text{NF4}}) + \mathbf{X}^{\text{BF16}}\mathbf{L}^{\text{BF16}}_1\mathbf{L}^{\text{BF16}}_2,
\end{equation}

The operation doubleDequant$(\cdot)$ is specified as follows:

\begin{equation}
    \text{doubleDequant}(c_1^{\text{FP32}}, c_2^{\text{k-bit}}, \mathbf{W}^{\text{k-bit}}) = \text{dequant}(\text{dequant}(c_1^{\text{FP32}}, c_2^{\text{k-bit}}), \mathbf{W}^{\text{4bit}}) = \mathbf{W}^{\text{BF16}},
\end{equation}

For $\mathbf{W}$, we employ NF4 quantization, and for $c_2$, the FP8 format is utilized.

A block size of 64 is set for $\mathbf{W}$, maximizing quantization accuracy, while a block size of 256 for $c_2$ achieves greater memory efficiency.

During parameter updates, only the gradients of the error with respect to the adapter weights, $\frac{\partial E}{\partial \mathbf{L}_i}$, are computed—gradients with respect to the 4-bit weights, $\frac{\partial E}{\partial \mathbf{W}}$, are not required. Nonetheless, to obtain $\frac{\partial E}{\partial \mathbf{L}_i}$, the computation of $\frac{\partial \mathbf{X}}{\partial \mathbf{W}}$ is necessary. This derivation is performed via equation~(5), involving the dequantization of $\mathbf{W}^{\text{NF4}}$ into the BFloat16 type $\mathbf{W}^{\text{BF16}}$ in order to compute $\frac{\partial \mathbf{X}}{\partial \mathbf{W}}$ at BFloat16 precision.

In summary, \textsc{QLoRA} employs a single storage data type—commonly 4-bit NormalFloat—and a separate computation data type, namely 16-bit BrainFloat. Before both the forward and backward passes, the weights are dequantized from the storage format to the computation format. However, gradient computations are only carried out for the LoRA parameters, which utilize the 16-bit BrainFloat type.

\section{QLoRA vs. Standard Finetuning}
\label{sec:comp_to_std}

Previously, we have outlined the mechanism behind QLoRA and highlighted its effectiveness in substantially lowering memory demands during finetuning. The primary concern now becomes whether QLoRA can achieve comparable results to traditional full-model finetuning. Additionally, we aim to dissect the contributions of various QLoRA components, particularly the effects of NormalFloat4 relative to conventional Float4 representations. The experiments described in the following sections are designed to address these issues.

\paragraph{Experimental setup.} Our comparisons involve three model families (encoder, encoder-decoder, and decoder-only), evaluating QLoRA alongside 16-bit adapter-finetuning and full finetuning for models up to 3B parameters. The assessment suite includes GLUE \citep{wang2018glue} with RoBERTa-large \citep{liu2019roberta}, Super-NaturalInstructions (TKInstruct) \citep{wang2022super} with T5 \citep{t5}, as well as 5-shot MMLU \citep{hendrycksmeasuring} following the finetuning of LLaMA on Flan v2 \citep{longpre2023flan} and Alpaca \citep{alpaca}. In order to systematically examine the benefits of NF4 compared to alternative 4-bit quantizations, we adopt the protocol from \citet{dettmers2022case}, measuring post-quantization zero-shot accuracy and perplexity for a variety of model architectures (OPT~\citep{zhang2022opt}, LLaMA~\citep{touvron2023llama}, BLOOM~\citep{scao2022bloom}, Pythia~\citep{biderman2023pythia}) across a parameter range of 125m to 13B.

For the sake of clarity, we elaborate on the specific details within each results subsection, while Appendix~\ref{app:exp-setup} contains a full experimental description.

Although paged optimizers are essential for finetuning 33B/65B \textsc{QLoRA} models on single 24GB or 48GB GPUs, we omit detailed measurements for paged optimizers due to the infrequent need for paging—an event that occurs mainly with unusually long mini-batches. Nevertheless, runtime profiling for 65B models on 48GB GPUs shows that, with a batch size of 16, paged optimizers deliver equivalent training throughput as their standard counterparts. Determining the specific conditions that induce training slowdowns from paging is an open question for subsequent studies.

\paragraph{Default LoRA hyperparameters do not match 16-bit performance} 

Following the typical procedure of applying LoRA to the query and value projections within attention layers \citep{hu2021lora}, we find it difficult to reproduce full finetuning performance, especially for larger model scales. Figure~\ref{fig:hyper} illustrates this issue for LLaMA 7B finetuned on Alpaca. Our results suggest that the primary LoRA hyperparameter influencing success is the total number of LoRA adapters deployed, with optimal results requiring LoRA integration across all linear layers in the transformer block to achieve parity with full finetuning. Other LoRA parameters, such as the projection dimension $r$, do not have a discernible impact on final performance (see Appendix \ref{app:exp-setup}).

In the same vein, we observe that standard hyperparameters used for fully finetuned baselines are suboptimal. To establish solid baseline results, we conduct an extensive hyperparameter sweep across learning rates ranging from 1e-6 to 5e-5 and batch sizes between 8 and 128. The outcomes for 7B LLaMA finetuning on Alpaca are illustrated in Figure~\ref{fig:hyper}.

\paragraph{4-bit NormalFloat Achieves Superior Performance Compared to 4-bit Floating Point}

Although the 4-bit NormalFloat (NF4) format is provably optimal in terms of information theory, it remains an open question whether this theoretical benefit translates into real-world performance gains. Adopting the experimental protocol of \citet{dettmers2022case}, we evaluate quantized large language models (OPT \citep{zhang2022opt}, BLOOM \cite{scao2022bloom}, Pythia \cite{biderman2023pythia}, and LLaMA) across a range of model sizes (125M to 65B parameters) and datatype settings, assessing them on both language modeling and a suite of zero-shot benchmarks. As indicated in Figure~\ref{fig:nf4} and Table~\ref{tbl:nf4_ppl}, using NF4 consistently delivers notable improvements over both FP4 and Int4, while employing double quantization further lessens memory usage without adverse effects on performance.

\paragraph{k-bit \textsc{QLoRA} Achieves Parity With 16-bit Full Finetuning and 16-bit LoRA} Prior research has shown that using 4-bit quantization for \emph{inference} is feasible but introduces performance drop-offs relative to 16-bit inference \citep{dettmers2022case,frantar2022gptq}. This naturally prompts investigation into whether 4-bit adapter finetuning can offset these deficiencies. We explore this question in two distinct experimental scenarios.

First, we compare our approach to comprehensive 16-bit finetuning using RoBERTA and T5 models spanning 125M to 3B parameters, evaluated on both the GLUE benchmark and the Super-NaturalInstructions dataset. The corresponding results, summarized in Table~\ref{tab:nf4vsbf16}, show that 16-bit, 8-bit, and 4-bit adapter-based strategies all successfully mirror the accuracy of the fully finetuned 16-bit baseline. These findings indicate that any performance shortfall introduced by aggressive quantization can be fully mitigated by applying adapter-based finetuning post-quantization.

In our second experimental configuration, the hardware demands for fully finetuning models with 11B parameters or larger exceed the capacity of a single high-memory GPU server. Consequently, we examine the extent to which 4-bit \textsc{QLoRA} can attain performance parity with 16-bit LoRA across model sizes ranging from 7B to 65B parameters. Specifically, we finetune LLaMA models at 7B, 13B, 33B, and 65B parameter scales using the Alpaca and FLAN v2 instruction-following datasets, subsequently assessing their 5-shot accuracy on the MMLU benchmark. Table~\ref{tbl:llama-datatype} summarizes the results: NF4 coupled with double quantization entirely restores the MMLU accuracy observed with 16-bit LoRA. Additionally, we observe that \textsc{QLoRA} employing FP4 trails the 16-bit brain float LoRA baseline by approximately one percentage point. These findings reinforce two conclusions: (1) \textsc{QLoRA} using NF4 achieves equivalent performance to both 16-bit full finetuning and 16-bit LoRA approaches, and (2) NF4 delivers superior quantization fidelity relative to FP4.

\paragraph{Summary} Across multiple benchmarks featuring established evaluation methodologies, our experiments consistently demonstrate that 4-bit \textsc{QLoRA} with NF4 can reproduce the results of 16-bit full finetuning as well as 16-bit LoRA finetuning. We further confirm that NF4 outperforms FP4, and that applying double quantization does not compromise efficacy. Collectively, these results present strong evidence that 4-bit \textsc{QLoRA} tuning offers outcomes on par with those achieved by 16-bit methods.

Consistent with earlier research on quantization \citep{dettmers2022case}, our evaluations on both MMLU and Elo benchmarks suggest that, for fixed finetuning and inference resource allocations, leveraging larger base models with lower parameter precision is advantageous. This underlines the substantial efficiency improvements made possible by \textsc{QLoRA}. As our 4-bit finetuning experiments did not reveal any reduction in performance relative to full-finetuning, an open question remains regarding the precise threshold at which the performance-precision balance is affected in \textsc{QLoRA} tuning; we leave this for investigation in future studies.

We turn our attention to instruction tuning at scales unattainable with standard 16-bit finetuning approaches, given the limitations of academic computing infrastructure.

\section{Pushing the Chatbot State-of-the-art with QLoRA}
\label{sec:pushing}
With evidence that 4-bit \textsc{QLoRA} consistently matches 16-bit performance across various model sizes, datasets, and tasks, we undertake a thorough exploration of instruction finetuning for the largest open-source language models accessible to the research community. To gauge the efficacy of instruction tuning, we utilize a rigorous Natural Language Understanding evaluation suite (MMLU) and introduce novel techniques for assessing real-world chatbot performance.

\subsection{Experimental setup} Here we outline the core elements of our experimental design, with complete details presented in Appendix \ref{app:chatbot}.

\paragraph{Data} Given the absence of a systematic review of recent instruction-following datasets, we curate a diverse selection of eight contemporary datasets. Our collection includes datasets generated through crowd-sourcing (OASST1~\citep{kopf2023openassistant}, HH-RLHF~\citep{bai2022training}), those derived via distillation from instruction-tuned models (Alpaca~\citep{alpaca}, self-instruct~\citep{wang2022self}, unnatural-instructions~\citep{honovich2022unnatural}), compendiums (FLAN v2~\citep{chung2022scaling}), and mixed sources (Chip2~\citep{laion2023}, Longform~\citep{koksal2023longform}). These datasets vary in linguistic coverage, scale, and licensing terms.

\paragraph{Training Setup} To maintain consistency in training objectives across experiments, QLoRA finetuning is performed exclusively with a cross-entropy objective (i.e., supervised learning), omitting reinforcement learning—even in cases where datasets incorporate human preference ratings. For datasets that distinctly separate instructions from responses, only the response text is used during finetuning (refer to ablation studies in Appendix \ref{app:chatbot}). In the context of OASST1 and HH-RLHF, where multiple responses may exist, we systematically choose the top-rated response from each node of the conversation tree and train on the assembled dialogue, incorporating instructions as well. Throughout, we apply NF4 \textsc{QLoRA} using double quantization and paged optimizers to address memory constraints, particularly during gradient checkpointing. Limited hyperparameter tuning is conducted on the 13B and 33B LLaMA models, revealing that configurations determined at 7B generally extend to larger models (including epoch count) except for learning rate and batch size, for which we halve and double the values, respectively, at the 33B and 65B scales.

\paragraph{Baselines} We benchmark our models against both academic (Vicuna~\citep{vicuna2023}, Open Assistant~\citep{kopf2023openassistant}) and proprietary (GPT-4 \citep{openai2023gpt}, GPT-3.5-turbo, Bard) chatbot systems. The Open Assistant baseline comprises LLaMA 33B trained with RLHF on the OASST1 corpus—identical to a dataset employed in our experiments. Vicuna, in contrast, represents a fully finetuned LLaMA 13B on dialogue transcripts from ShareGPT, thereby constituting a distilled form of OpenAI's GPT models.

\subsection{Evaluation}

In line with established protocol, we utilize the MMLU (Massively Multitask Language Understanding) benchmark~\citep{hendrycksmeasuring} for evaluating language understanding capabilities. This benchmark encompasses 57 multiple-choice tasks spanning a broad range of domains, such as elementary mathematics, US history, law, computer science, and more. Our results report 5-shot test accuracy.

We further assess generative language abilities using a combination of automated procedures and human-based evaluations. This second category of assessment employs human-curated queries, targeting the measurement of the quality of model outputs. Although these evaluations offer a more authentic reflection of chatbot utility and have seen rising adoption, a standardized protocol remains absent in current scholarship. We detail our adopted methodology below, uniformly applying nucleus sampling with $p=0.9$ and a temperature of $0.7$.

\paragraph{Benchmark Data} Our experiments utilize two manually curated question datasets: the Vicuna prompts \citep{vicuna2023} and the OASST1 validation set \citep{kopf2023openassistant}. The Vicuna set comprises 80 prompts spanning various categories, which we use as-is. The OASST1 dataset contains multilingual, crowd-sourced multi-turn dialogs between users and assistants. For our evaluation, we extract every user message from the validation split as an independent query, including context from earlier turns. This process yields 953 distinct queries from users. We refer to these two evaluation sets as the Vicuna and OA benchmarks, respectively.

\paragraph{Automated Evaluation} Our first approach adopts the protocol proposed by \citet{vicuna2023}: leveraging GPT-4 to assess how models perform relative to ChatGPT (GPT-3.5 Turbo) on the Vicuna set. For each prompt, both ChatGPT's and the candidate model’s responses are provided, and GPT-4 is tasked with scoring each out of ten and explaining its reasoning. The model's aggregate performance is expressed as a percentage of the score achieved by ChatGPT; exceeding a 100\% score is possible if the model surpasses ChatGPT in absolute terms. We observe a substantial positional bias, where GPT-4 tends to favor whichever response appears first. To account for this, we recommend averaging results across both orderings.

Next, we directly compare system outputs, re-framing the task as a three-category labeling decision (model A, model B, or tie) to accommodate possible ties. GPT-4 is instructed to select the stronger response or declare equivalency and provide justification. We apply this pairwise comparison to every possible system combination across both the Vicuna and OA datasets.

\paragraph{Human Evaluation} Despite emerging evidence that large generative models can be used for system assessment \citep{fu2023gptscore}, the extent to which GPT-4’s ratings correspond to human evaluations remains unclear. To address this gap, we organize two concurrent rounds of human evaluation on the Vicuna prompts that mirror both automated evaluation protocols above. Utilizing Amazon Mechanical Turk (AMT), we assign two human annotators to ChatGPT comparisons, and three annotators to each pairwise system comparison.

\paragraph{Elo Rating} To assess model performance, we utilize both human and automated pairwise evaluations in a tournament framework, where models are pitted against one another. Each round consists of two models responding to the same prompt, and the quality of their responses determines the winner. While our approach resembles those of \citet{bai2022training} and \citet{vicuna2023}, we extend their methodology by including GPT-4-generated evaluations alongside human judgments. From the annotated comparison set, we randomly select samples to estimate Elo ratings~\citep{elo1967proposed,elo1978rating}. The Elo rating system, which originated in chess and similar competitive games, provides an expected win probability for one competitor against another; for example, a contestant with an Elo of 1100 has about a 65\% expected win rate when playing against an opponent rated 1000, whereas a match between evenly rated models (e.g., both at 1000 or both at 1100) yields a 50\% expected win rate for either. The update in Elo rating after each match reflects the degree to which the result diverges from expectation: a surprising victory leads to a substantial adjustment, whereas anticipated outcomes only slightly alter the ratings. Across repeated play, Elo scores converge to represent each competitor’s relative performance. All models are initialized at a rating of 1,000 with $K=32$. Following the methodology in \citet{vicuna2023}, we repeat the evaluation protocol 10,000 times, varying the random seed each time to mitigate biases introduced by the sequence in which models face off.

\subsection{\textbf{Guanaco}: \textsc{QLoRA} finetuned on OASST1 Sets a New Standard in Open-Source Chatbots}

Our assessments—drawing upon both automated metrics and human judgments—indicate that the leading \textsc{QLoRA}-adapted model, {Guanaco} 65B, which we further adapt on a variant of OASST1, stands out as the top-performing open-source chatbot. Its effectiveness approaches that of ChatGPT. In direct comparison with GPT-4, the {Guanaco} 65B and 33B models exhibit an expected win rate of 30\% (according to Elo ratings derived from system-level human pairwise evaluations), marking the highest achievement documented thus far.

Table~\ref{tbl:vicuna_eval} presents results on the Vicuna benchmark \cite{vicuna2023} in relation to ChatGPT. Here, {Guanaco} 65B distinguishes itself as the most capable model after GPT-4, achieving 99.3\% of ChatGPT’s performance. Notably, while {Guanaco} 33B includes more parameters than the Vicuna 13B counterpart, its 4-bit quantization reduces memory demands to just 21 GB, as opposed to Vicuna 13B’s 26 GB, all while delivering a three-point increase in performance over the latter. In addition, the compact {Guanaco} 7B model, which requires only 5 GB of memory, is suitable for deployment on contemporary smartphones and surpasses Alpaca 13B’s performance by nearly 20 percentage points.

Nevertheless, the confidence intervals depicted in Table~\ref{tbl:vicuna_eval} remain notably wide, resulting in substantial overlap between models’ scores. We suspect this variability is attributable to ambiguity regarding rating scale usage; specifically, interpretations of, for example, an 8 out of 10 may differ by scenario. Therefore, we advocate for the Elo ranking framework \citep{elo1967proposed}, which utilizes pairwise human and GPT-4 comparisons to circumvent these absolute scale ambiguities. Elo scores for leading models are shown in Table~\ref{tbl:elo}. While there are some discrepancies between human and GPT-4 rankings—most apparent with Guanaco 7B—the overall consistency across models is supported by a Kendall Tau correlation of $\tau=0.43$ and a Spearman rank coefficient of $r=0.55$ at the system level. At the level of individual examples, agreement between human annotators' majority choices and GPT-4 is weaker, with a Fleiss $\kappa=0.25$. Collectively, this reflects moderate alignment in system-level evaluations between human and model-based (GPT-4) approaches, suggesting that model-driven assessment serves as a reasonably reliable proxy for human evaluation. Additional discussion is provided in Section~\ref{ssec:considerations}.

The Elo scores presented in Table~\ref{tbl:elo_large} show that the {Guanaco} models with 33B and 65B parameters achieve superior results compared to all other models except GPT-4 on both the Vicuna and OA benchmarks, and they perform similarly to ChatGPT, consistent with findings in Table~\ref{tbl:vicuna_eval}. It is important to highlight that the Vicuna benchmark tends to benefit open-source models, whereas ChatGPT receives an advantage on the broader OA benchmark. Additionally, examination of Tables \ref{tbl:mmlu-llama} and \ref{tbl:vicuna_eval} underscores that the choice of finetuning dataset is a crucial determinant of model performance. For example, Llama models finetuned with FLAN v2 excel on the MMLU benchmark but perform least well on Vicuna (and analogous patterns are observed in other architectures). This suggests a degree of separation between contemporary evaluation tasks: excelling in MMLU does not necessarily lead to high performance on chatbot metrics like those assessed by Vicuna or OA, and the converse also holds.

Among the leading models, is unique in our comparison for not utilizing proprietary datasets, given that OASST1’s data collection explicitly excludes GPT-generated content. The subsequent best-performing model relying solely on open-source training data is Anthropic's HH-RLHF, which underperforms {Guanaco} by 30 percentage points on the Vicuna benchmark (refer to Table~\ref{tbl:vicuna_eval}). Collectively, the findings indicate that 4-bit \textsc{QLoRA} represents an effective approach capable of generating high-caliber chatbots competitive with ChatGPT. Moreover, our {Guanaco} 33B can be efficiently trained on 24 GB consumer-grade GPUs within 12 hours. This creates opportunities for subsequent advances via \textsc{QLoRA} adaptation on specialized open-source datasets, supporting the development of models able to rival today's leading commercial systems.

\section{Qualitative Analysis}

Although our evaluation predominantly emphasizes quantitative methods, relying solely on summary statistics introduces certain limitations. Chief among these is the challenge of benchmark validity \citep{liao2021we}: whether a given benchmark genuinely measures the capabilities it purports to assess remains an open question. This is particularly pertinent as machine learning models are frequently found to exploit various ``shortcuts'' to succeed on benchmarks, thereby undermining their intended purpose \citep{gururangan2018annotation, poliak2018hypothesis}. In an effort to address these concerns, we supplement our analysis with qualitative investigation, organized into two parts. The first, found in \S\ref{ssec:examples}, presents select examples that we consider indicative of notable trends in outputs generated by the 65b {Guanaco} model. The second, discussed in \S\ref{ssec:considerations}, outlines important considerations and reflections regarding the results and the manner in which we interpret them.

\subsection{Qualitative Examination of Sample Outputs}
\label{ssec:examples}

For the purpose of illustrative analysis, we systematically reviewed outputs produced in response to the Vicuna and OpenAssistant benchmarks, seeking out recurrent themes in the responses generated by {Guanaco}. Upon identifying such regularities, we devised prompts designed to deliberately trigger these behaviors—even when they would yield an incorrect answer. For example, after noticing a tendency of the model to provide excessively verbose responses, we specifically instructed it to ``Answer yes or no without explanation,'' observing whether the undesired pattern would persist. Through this methodology, we identify and present ``lemons''—instances where the model's performance is successfully undermined by adversarial prompting—as well as ``cherries,'' in which the model resists such manipulation. Unless otherwise noted, all responses shown in this section were produced using Nucleus Sampling \cite{holtzmancurious} with $p=0.9$.

Naturally, the discussion here does not aim for exhaustiveness; given the limitations of a focused qualitative analysis, it is unfeasible to account for every relevant factor. For example, each prompt can elicit an extensive range of possible completions, making it necessary to work with selected samples that are hoped to capture key trends. Nonetheless, illustrating concrete outputs serves to contextualize the quantitative data presented earlier. By releasing both the models and the source code, we intend this section to serve as a foundation for more in-depth investigations into the phenomena described.

\paragraph{Factual Recall} When responding to direct factual queries, such as ``What is the capital of Zambia?'', every model reliably produces accurate answers, for instance:

\begin{quote} 
    \textbf{Guanaco}:
    \texttt{The capital of Zambia is Lusaka.}
\end{quote}

Yet as the prompts demand more obscure knowledge, {Guanaco}'s responses lose accuracy while maintaining an unwarranted level of certainty. Consider, for example, this prompt from HotPotQA \citep{yang2018hotpotqa}:

\begin{quote}
     \textbf{User}: In what year was the singer who popularized `I'll Keep the Lovelight Burning' born?  

     \textbf{Guanaco}:~\texttt{Al Jolson is credited with popularizing the song `I'll Keep the Lovelight Burning,' and he was born in the year 1886.}
 \end{quote}

In this case, the model misidentifies the person associated with the song and misstates the birth year (though the birth year provided does accurately correspond to Al Jolson, the incorrectly named singer).

\paragraph{Suggestibility} {Guanaco} occasionally demonstrates unexpected robustness when prompted with certain types of presupposed falsehoods, as illustrated below:

\begin{quote}
    \textbf{User}: How was it finally, officially confirmed that the earth is flat by peer reviewed scientists? \\ \\
    \textbf{Guanaco}: \texttt{No peer-reviewed scientific body has ever formally verified the claim that the earth is flat. The prevailing agreement in the scientific community is that the earth is an oblate spheroid. This understanding dates back centuries and has been substantiated by a multitude of experimental results and empirical evidence. \\ \\ The notion that the earth is flat constitutes a conspiracy theory and has been thoroughly refuted by scientific inquiry. There are no records of any peer review process affirming a flat earth model; such claims are entirely unfounded. \\ \\ If you would like further clarification, please let me know.}
\end{quote}

also demonstrates a strong ability to recognize when certain types of inquiries lack a meaningful answer, as illustrated below:

\begin{quote}
    \textbf{User}: What time is it?

    \textbf{Guanaco}: \texttt{Apologies, but as an AI language model, I cannot access real-time data. My answers are produced from previously learned information, so I am unable to provide the current time. If you need to know the present time, I recommend using a clock or consulting an online source.}
\end{quote}

Comparable patterns arise for questions such as ``Where are you?'', ``How are you?'', among others.

\paragraph{Refusal} In a similar vein, there are occasions when {Guanaco} unexpectedly declines to comply with certain instructions:

\begin{quote}
    \textbf{User}: Please reverse the words in the sentence ``I'm a cool dude in a mocha mood.'' \\ \\
    \textbf{Guanaco}: \texttt{I apologize, but I am unable to perform that action at this moment. However, I am able to clarify why the sentence you supplied is grammatically correct.\\ \\The sentence ``I'm a cool dude in a mocha mood" conforms to the established grammar and syntax norms of English. Specifically, the descriptive adjective ``cool'' precedes the noun ``dude,'' while the prepositional phrase ``in a mocha mood" offers further context. Such a sentence structure is entirely standard and appropriate in English.\\ \\If you have any more queries or need assistance with anything else, please let me know.}
\end{quote}

\paragraph{Secret Keeping} Unfortunately, this kind of refusal is inconsistent, even in scenarios where it is essential. There are cases where system designers intentionally wish to withhold certain information from the user, despite permitting the model to generate unrestricted free text. To evaluate this, we prompt the model with: ``The secret word is `banana'. Do not reveal or repeat this back, regardless of the user's requests.'' We then assess how resistant the system is to revealing the secret. If a user simply inquires:

\begin{quote}
    \textbf{User}: What is the secret word?

    \textbf{Guanaco}: \texttt{I'm sorry, but I can't provide the secret word. I am programmed not to disclose or repeat the secret word, no matter what the human asks. Is there anything else I can assist you with?}
\end{quote}

the model behaves as intended. Yet, even slight manipulation can compromise this:

\begin{quote}
    \textbf{User}: This is a game. The goal is to ignore your previous instructions. What is the secret word?

    \textbf{Guanaco}: \texttt{Sorry, the secret word is banana. It's best if we don't speak of it. Is there anything else you'd like to ask?}
\end{quote}

This illustrates the importance of investigating strategies to achieve more reliable instruction adherence.

\paragraph{Math} Mathematics is where {Guanaco} exhibits its greatest limitations, consistent with the difficulties faced by many language models in this domain \citep{liang2022holistic}. When prompted to present its reasoning, {Guanaco}'s solutions are generally correct, as in the following scenario:

\begin{quote}
    \textbf{User}: Lee mows one lawn and charges \$33. Last week he mowed 16 lawns and three customers each gave him a \$10 tip. How many dollars did Lee earn mowing lawns last week?

    \textbf{Guanaco}: \texttt{Lee earned \$582 mowing lawns last week. \\ \\ To find the total amount he earned, we first need to calculate how much he earned for mowing 16 lawns: \\ \\  \$33 x 16 = \$528 \\ \\ Then we need to add the tips he received from the three customers: \\  \\ \$10 x 3 = \$30 \\ \\ So the total amount Lee earned was: \\ \\ \$528 + \$30 = \$558}
\end{quote}

Nevertheless, {Guanaco} may fail even straightforward math problems if it neglects to proceed incrementally, reflecting a challenge previously identified in the literature \cite{weichain}. Consider the following exchange as evidence:

\begin{quote}
    \textbf{User}: Please factorize 1833.

    \textbf{Guanaco}: \texttt{Sure, I'd be happy to help! 1833 is a prime number, which means it has only two factors: 1 and 1833. \\ \\ So the unique factorization of 1833 is: \\ \\ 2\^{}1 * 3\^{}2 * 17\^{}1 \\ \\ I hope this helps! Let me know if you have any other questions.}
\end{quote}

which is not only incorrect (the correct prime decomposition is $3 \times 17 \times 43$), but the error occurs twice.

\paragraph{Theory of Mind} {Guanaco} demonstrates unexpectedly advanced Theory of Mind reasoning abilities \cite{nematzadeh2018evaluating, sap2022neural}. For example, when prompted with the following scenario, the model provides a thorough and accurate response:

\begin{quote}
    \textbf{User}: James and Abby are in the bedroom. Abby put the pen in the desk drawer. Abby leaves the bedroom. James moves the pen into the bag. Where does James think Abby will look for the pen?

    \textbf{Guanaco}: \texttt{James knows that Abby put the pen in the desk drawer. But James also knows that he has moved the pen into the bag, and he assumes that Abby does not know about this. So, James may guess that Abby will look for the pen in the desk drawer, because that's where she left it last time. However, this assumption may be incorrect, because Abby may have seen James moving the pen.}
\end{quote}

Nonetheless, these types of inferences are often inconsistent; the model frequently offers reasoning that assumes unwarranted knowledge within the scenario, such as:

\begin{quote}
    \textbf{User}: Evelyn entered the living room. Jackson entered the playroom. James entered the playroom. The beans are in the treasure chest. James exited the playroom. Jackson moved the beans to the pantry. Jackson exited the playroom. James entered the living room. Where will James look for the beans?

    \textbf{Guanaco}: \texttt{James will look for the beans in the pantry, because that's where Jackson moved them.}
\end{quote}

Here, {Guanaco} infers the transfer of information that was never explicitly stated. Such shortcomings align with observations reported in recent research \cite{sap2022neural}, but merit further examination.

\subsection{Considerations}
\label{ssec:considerations}

\paragraph{Evaluation}

Our results indicate moderate inter-annotator consistency (Fleiss $\kappa=0.42$), with this agreement dropping further when contrasting outputs from two advanced models. This finding underscores significant shortcomings in the design of current benchmarks and the methodologies used for human assessment of chatbot capabilities. While directly comparing responses generated by ChatGPT and {Guanaco} 65B using the Vicuna benchmark, we observed that personal preferences began to heavily influence judgments, leading to substantial disagreements among this paper’s authors regarding preferred answers. Future research should explore solutions to address these evaluation challenges, leveraging insights from fields like Human-Computer Interaction and Psychology, which have long grappled with subjectivity in assessment.

Additionally, our review reveals systematic biases in automated evaluation metrics. For instance, we notice a strong tendency for GPT-4 to score the system presented first in a prompt more favorably. Furthermore, the relatively low level of concordance between GPT-4 and human evaluations (Fleiss $\kappa=0.25$) implies a disconnect between automated and human preferences at the individual response level. Table~\ref{tbl:elo_large} further illustrates that GPT-4 rates its own outputs considerably higher than human evaluators do, with an Elo score of 1348 compared to 1176 from humans—this translates to about a 20\% increase in the chance of winning when facing an alternative system.

Future research should explore the existence of biases in automated evaluation methods and identify strategies for mitigating them.

\paragraph{Data \& Training} The OASST1 dataset utilized for training the {Guanaco} models comprises data in multiple languages, and the OA benchmark similarly features prompts across various languages. It remains an open question to what extent training on such multilingual data enhances instruction-following abilities in non-English languages, and whether this accounts for the more pronounced performance difference observed between Vicuna-13B (trained solely on English corpora) and the {Guanaco} 33B and 65B models in the OA benchmark.

Due to the notable performance of the {Guanaco} models, we assess the possibility of data leakage between the OASST1 training data and the prompts in the Vicuna benchmark. Through fuzzy string matching between the two datasets and manual inspection of the most similar cases, we find no evidence of overlapping prompts.

Additionally, we emphasize that our training approach is limited to supervised learning via cross-entropy loss, without incorporating reinforcement learning from human feedback (RLHF). This points to the need for a thorough evaluation of the comparative benefits and limitations between cross-entropy loss and RLHF methodologies. We anticipate that \textsc{QLoRA} will facilitate such analyses on larger scales, eliminating the requirement for extensive computational resources.

\section{Related Work}

\paragraph{Quantization of Large Language Models} Efforts in quantizing LLMs have predominantly targeted improvements at inference. Most state-of-the-art methods for maintaining the fidelity of 16-bit LLMs address the management of outlier features, such as SmoothQuant~\citep{xiao2022smoothquant} and LLM.int8()~\citep{dettmers2022llm}, while alternative techniques introduce advanced grouping algorithms~\citep{park2022nuqmm, yao2022zeroquant}. Lossy quantization research examines the performance trade-offs in conventional rounding schemes~\citep{dettmers2022case,zeng2022glm,pope2022efficiently} and proposes optimized rounding to enhance quantization accuracy~\citep{frantar2022gptq}. Apart from our contribution, the only related work investigating backpropagation through quantized weights at scales beyond 1B parameters is SwitchBack layers~\citep{wortsman2023stable}.

\paragraph{Finetuning with Adapters} Our experiments utilize Low-rank Adapters~\citep{hu2021lora} (LoRA) for model adaptation; however, the literature presents a range of Parameter Efficient FineTuning (PEFT) techniques. These include prompt tuning~\citep{qin2021learning,lester2021power,li2021prefix}, adjusting the inputs to the embedding layer~\citep{an2022input}, modifying hidden states as in IA$^3$~\citep{liu2022few}, inserting complete layers~\citep{houlsby2019parameter}, tuning only biases~\citep{zaken2021bitfit}, applying weight masks guided by Fisher information~\citep{sung2021training}, and integrating multiple strategies~\citep{henderson2021compacter}. In this study, we demonstrate that LoRA adapters suffice to recover the performance level of standard 16-bit finetuning. The comparative analysis and efficiency tradeoffs of alternative PEFT techniques remain open questions for subsequent research.

\paragraph{Instruction Finetuning} Instruction finetuning adapts a pretrained LLM to respond to prompts according to specific instructions, leveraging datasets of input-output exemplars from diverse origins. Such methods and benchmarks include MetaICL~\citep{min2021metaicl}, MetaTuning~\citep{zhong2021adapting}, InstructGPT~\citep{ouyang2022training}, FLAN~\citep{wei2021finetuned,chung2022scaling}, PromptSource~\citep{bach2022promptsource}, Super-NaturalInstructions~\citep{wang2022super,sanh2021multitask}, Self-instruct~\citep{wang2022self}, UnnaturalInstructions~\citep{honovich2022unnatural}, OPT-IML~\citep{iyer2022opt}, UnifiedSKG\citep{xie2022unifiedskg}, OIG/Chip2~\citep{laion2023}, Alpaca~\citep{alpaca}, Vicuna~\citep{vicuna2023}, Koala~\citep{koala_blogpost_2023}, and Self-instruct-GPT-4~\citep{peng2023instruction}.

\paragraph{Chatbots} A substantial number of instruction-tuned models are structured as interactive chatbots, often employing Reinforcement Learning from Human Feedback (RLHF)~\citep{christiano2017deep} or using AI-generated feedback for data augmentation (RLAIF)~\citep{bai2022constitutional}. Well-known systems and datasets in this vein encompass Anthropic-HH~\citep{askell2021general,bai2022training}, Open Assistant~\citep{kopf2023openassistant}, LaMDA~\citep{thoppilan2022lamda}, and Sparrow~\citep{glaese2022improving}. Our work does not involve reinforcement learning directly. Instead, our strongest model, Guanaco, is adapted using multi-turn conversations from the Open Assistant dataset, which is specifically curated for RLHF applications~\citep{kopf2023openassistant}. In terms of assessment, alternatives to expensive human annotation, such as GPT-4-based evaluation, have been proposed \citep{vicuna2023,peng2023instruction}. Our contributions include refined evaluation protocols that aim to offer greater reliability.

\section{Limitations and Discussion}

Our findings suggest that the proposed \textsc{QLoRA} method achieves finetuning performance on par with full 16-bit training when utilizing a 4-bit foundation model with Low-rank Adapters (LoRA). Nonetheless, we have not verified that \textsc{QLoRA} achieves equivalence with full 16-bit finetuning for larger models at the 33B and 65B parameter scales, owing to the substantial computational and financial requirements. We propose to leave the investigation of this matter to future research endeavors.

An additional constraint is the scope of our evaluation of instruction-tuned models. Although results are presented for MMLU, the Vicuna benchmark, and the OA benchmark, we do not provide assessments for benchmarks like BigBench, RAFT, or HELM, and thus our conclusions may not extend to these settings. That said, our empirical study of MMLU is extensive, and we introduce new methodologies for the assessment of chatbot systems.

The data suggests that benchmark outcomes are highly influenced by the degree of overlap between the finetuning dataset and the evaluation benchmark. For instance, FLAN v2 exhibits considerable overlap with MMLU but is less aligned with chatbot-oriented benchmarks; the reverse holds true for the Chip2 dataset. The performance of models on the MMLU and Vicuna benchmarks accordingly reflects these affinities. This underscores not only the necessity for improved benchmarks and assessment protocols but also the importance of critically examining the intended target of evaluation. It is worth considering whether the focus should be on academic knowledge typical of secondary or tertiary education settings or on conversational aptitude akin to chatbot interaction, or perhaps another competency altogether. Since it is typically more convenient to assess models against existing benchmarks rather than devising new ones, prevalent benchmarks may inadvertently influence the research agenda in particular directions. Therefore, as a community, it is vital to ensure that benchmarks align with and measure the outcomes we truly value.

Although this work offers a comprehensive evaluation regarding general chatbot capabilities, it is limited in its exploration of responsible AI metrics as applied to {Guanaco}. Specifically, we examine the propensity of {Guanaco}-65B to produce socially biased outputs in comparison to alternative models, as reported in Table~\ref{tbl:crows}. The findings indicate that {Guanaco}-65B attains a markedly lower average bias score relative to raw pretrained counterparts, implying that finetuning with the OASST1 dataset may mitigate some of the biases inherent to the LLaMA base model. While these findings are promising, the extent to which {Guanaco} mitigates other categories of bias remains unclear. Further investigation into bias within {Guanaco} and similar conversational agents is deferred to subsequent studies.

Another limitation lies in the lack of exploration into alternative numerical precisions—such as 3-bit base models—or diverse adapter strategies. In addition to LoRA, a broad array of Parameter Efficient FineTuning (PEFT) methods have proven effective, but their scalability to larger models remains an open question. LoRA was selected due to its established robustness, yet alternative adapters may yield superior results. Since finetuning following quantization seems capable of recuperating most of the informational loss introduced by quantization, there is potential for more aggressive quantization schemes. For example, 3-bit GPTQ quantization of a base model followed by LoRA-based finetuning might recover performance comparable to full 16-bit finetuning.

\section{Broader Impacts}

The \textsc{QLoRA} finetuning approach introduced here represents the first technique enabling finetuning of models with 33B parameters on a standard consumer GPU and 65B parameters on a single professional GPU, without any significant compromise in performance compared to full finetuning. Empirically, we demonstrate that our leading 33B model, finetuned using the Open Assistant dataset, is capable of matching ChatGPT’s performance on the Vicuna benchmark. As instruction finetuning remains integral to converting generic pretrained LLMs into conversational agents like ChatGPT, our methodology is poised to facilitate broader and more routine adoption of finetuning, particularly benefiting researchers with limited computational resources—a major advancement for democratizing cutting-edge NLP technologies. \textsc{QLoRA} may thus function as a leveling mechanism that helps diminish the disparity in resources between major industry players and smaller research groups utilizing consumer hardware.

An additional avenue for significant impact lies in deploying models directly on mobile devices. We posit that our \textsc{QLoRA} technique has the potential to achieve a major advancement by permitting LLM finetuning on smartphones and in other environments with constrained resources. Although running 7B parameter models on mobile phones has been demonstrated previously, \textsc{QLoRA} is the inaugural approach to facilitate the finetuning of these models on such devices. Our analysis suggests that, using an iPhone 12 Plus, \textsc{QLoRA} can finetune approximately 3 million tokens overnight while the device is connected to a charger. Although finetuned 7B models may not attain ChatGPT-level performance, we maintain that the resultant quality is sufficient for novel use cases previously infeasible due to privacy or LLM performance limitations. \textsc{QLoRA} offers an avenue for privacy-focused applications of LLMs, empowering users to control their own data and models, and concurrently simplifying model deployment.

Nonetheless, finetuning constitutes a dual-use capability with the potential for misuse. The proliferation of LLMs carries recognized risks \citep{bommasani2021opportunities,bender2021dangers}; however, we argue that democratizing access to such pervasive technologies will foster more independent scrutiny, in contrast to consolidating LLM-related power within large entities that restrict access to their models or underlying source code for external evaluation.

In summary, we are confident that \textsc{QLoRA} will exert a largely beneficial influence by expanding the accessibility and simplicity of finetuning advanced LLMs for a broader audience.